\documentclass[11pt,a4paper]{article}
\usepackage[utf8]{inputenc}
\usepackage[T1]{fontenc}
\usepackage[french]{babel}
\usepackage{amsmath,amssymb}
\usepackage[margin=2.4cm]{geometry}
\usepackage{booktabs}
\usepackage{multirow}
\usepackage{natbib}
\usepackage{hyperref}
\hypersetup{colorlinks=true,linkcolor=blue,citecolor=blue,urlcolor=blue}

\title{\bf Les critères pré-enregistrés sont des spécifications :\\
onze défaillances observées dans une campagne de cosmologie,\\
cartographiées sur des pathologies connues}
\author{Édouard Lantenois\thanks{\texttt{github.com/Dantenos}}}
\date{Août 2026}

\begin{document}
\maketitle

\begin{abstract}
\noindent
Nous rapportons onze façons dont un critère de succès gelé par empreinte a échoué au cours
d'une campagne intensive de cosmologie --- 89 critères gelés, 201 entrées de registre, 81
attaques arbitrées en deux jours --- et nous cartographions chacune sur une pathologie déjà
nommée hors de notre domaine. \textbf{Aucune des onze n'est un mode de défaillance nouveau.}
L'une d'elles --- une validation passant à $10^{-4}$ en valeur relative alors que la grandeur
qu'elle convertit n'existe pas --- est la \emph{vacuité}, nommée par \citet{beer1997} en 1997,
dont l'exemple canonique (« toute requête est finalement acquittée », satisfaite par un
système qui n'émet jamais de requête) est le nôtre avec une tolérance flottante à la place
d'un opérateur temporel. Une autre est l'incohérence code--commentaire ; une autre un
paramètre au bord du domaine physique \citep{feldman1998} ; une autre une spécification
insatisfiable ; une autre un test instable sur égalité flottante. Notre revendication est
donc une \emph{cartographie}, non une découverte : les critères d'acceptation
pré-enregistrés sont des spécifications, et ils héritent de toute la pathologie connue des
spécifications. Nous soutenons que cela mérite d'être énoncé parce que les taxonomies
existantes de défaillance du pré-enregistrement \citep{yamada2018,bakker2020} sont
\emph{comportementales} --- elles classent ce que fait un chercheur --- alors que les
défaillances observées ici sont \emph{structurelles} : elles classent ce que le critère
\emph{est}, et chacune s'est produite sans la moindre malhonnêteté. Nous donnons le décompte
des instances, la classe de chacune, et quel outillage existant l'aurait attrapée. Nous
rapportons un invariant de protocole pour lequel nous n'avons trouvé aucune antériorité ---
le bloc des critères est immuable tandis que le corps du script reste librement modifiable
--- et nous renonçons explicitement à revendiquer le reste du protocole, qui est une
intégration d'enregistrement à la OSF, d'intégration continue \citep{bj2017} et de
programmation par contrat. Enfin nous rapportons la seule défaillance qui n'est pas celle
d'un critère : après trois conceptions gelées bloquées sur la même question, une quatrième
serait passée, et l'écrire aurait été du p-hacking d'un cran au-dessus. Le gel par empreinte
protège chaque étude ; il ne protège pas la série.
\end{abstract}

\section{Le protocole, rapporté comme note de conception}
\label{sec:protocole}

Nous énonçons le mécanisme brièvement, et nous ne le revendiquons pas. L'unité est un script.
Son docstring s'ouvre sur un bloc de critères de succès : les validations à passer avant toute
publication, les critères eux-mêmes --- exigés \emph{exhaustifs} et \emph{exclusifs}, de sorte
que chaque issue corresponde à exactement un verdict écrit --- et une déclaration de ce que
l'étude ne revendiquera pas. \texttt{freeze} extrait le docstring par AST et enregistre son
SHA-256 ; \texttt{verify} fait échouer l'intégration continue si l'empreinte a bougé ; les
amendements sont permis mais consignés dans un journal public avec justification obligatoire.

Chaque composant est disponible sur étagère. Les critères dans l'artefact exécutable, c'est la
programmation par contrat, et en Python spécifiquement c'est \texttt{doctest}. Geler un plan
d'analyse par empreinte, c'est l'enregistrement OSF, et sous forme cryptographique OriginStamp
\citep{originstamp} et CryptSubmit \citep{cryptsubmit} ; une recherche dans OpenAlex rend de
l'ordre d'une centaine de dépôts en physique et cosmologie le faisant déjà, dont plusieurs de
2026 combinant un manifeste SHA-256 et une attestation OpenTimestamps. L'application par
intégration continue, c'est l'analyse continue \citep{bj2017} et la convention Popper
\citep{popper}. Un journal d'amendement public avec justification, c'est le mécanisme de
retrait d'OSF, et le signalement des déviations a sa propre littérature de guidage
\citep{willroth2024}. Fixer les critères d'acceptation avant de regarder, c'est l'analyse
aveugle, standard en physique des particules depuis vingt ans \citep{klein2005,maccoun2015} et
en cosmologie depuis DES \citep{muir2020}.

\paragraph{Le seul invariant que nous n'ayons pas trouvé ailleurs.} L'antériorité gèle
l'\emph{artefact} --- le plan, le code, le manifeste. Nous ne gelons que le \emph{prédicat
d'acceptation}, et laissons le corps librement modifiable. Cette asymétrie est ce qui rend
correction et pré-enregistrement compatibles : un script peut être débogué, reparamétré et
relancé sans toucher à ce qui compterait comme un succès. Sur 89 critères gelés, le canal
d'amendement a servi zéro fois --- non parce que les critères étaient toujours justes, mais
parce qu'avoir tort était survivable sans les modifier. Nous offrons cette asymétrie comme
seule contribution du protocole et ne défendrions rien de plus.

\section{La thèse : les critères sont des spécifications}
\label{sec:these}

Un critère pré-enregistré est un prédicat qu'un calcul doit satisfaire pour qu'une
affirmation compte comme établie. C'est la définition d'une spécification. La conséquence est
immédiate et, autant que nous puissions le vérifier, jamais énoncée : \textbf{les critères
scientifiques pré-enregistrés héritent de toute la pathologie connue des spécifications} ---
satisfaction vacue, insatisfiabilité, dérive vis-à-vis de l'implémentation, sensibilité
insuffisante --- ainsi que de la littérature de détection qui accompagne chacune.

Nous y sommes arrivés par le chemin long. Nous avons écrit onze critères qui ont échoué, nous
les avons catalogués comme s'ils étaient neufs, et nous avons découvert en vérifiant que
chacun portait déjà un nom.

\section{Onze instances observées, et la classe de chacune}
\label{sec:mapping}

\begin{table}[t]
\centering
\small
\begin{tabular}{p{4.9cm}p{3.6cm}p{5.9cm}}
\toprule
instance observée & classe nommée & d'où vient le nom \\
\midrule
Critère gelé sur une grandeur que les données ne peuvent pas contraindre ; satisfait par
toute issue. &
\textbf{Non-identifiabilité} & théorie de l'estimation, standard \\
\addlinespace
Une borne de qualité d'ajustement $\chi^2/(N+17)$ gelée pour un jeu à 22 intervalles ; le
décalage venait d'une autre configuration de données. &
\textbf{Mésusage du $\chi^2$ réduit} & \citet{andrae2010} \\
\addlinespace
Une validation comparant le modèle en un point \emph{fixé} à une référence prise à son
\emph{optimum} ; l'écart mesuré était l'optimisation. &
\multirow{3}{3.4cm}{\textbf{Comparant faible ou inadapté} --- ces deux-là n'en font qu'une, et
un référé le dirait} &
\multirow{3}{5.7cm}{\citet{dacrema2019} ; le domaine « biais de comparant » des instruments
cliniques de risque de biais} \\
Une validation jouée contre un comparant qui réintroduisait un terme que le modèle avait
retiré. & & \\
 & & \\
\addlinespace
Une contrainte serrée dont le minimum était sur le bord du domaine accessible ; une branche
était fermée par l'implémentation, donc un résultat unilatéral était rapporté bilatéral. &
\textbf{Paramètre au bord physique} & \citet{feldman1998}, dont c'est le sujet entier :
intervalles uni- contre bilatéraux au bord \\
\addlinespace
Deux exigences gelées dans un même bloc --- une grille de 41 points et une tolérance de
$\pm0{,}3$ --- qui ne pouvaient pas être satisfaites ensemble. L'étude ne pouvait pas réussir
par construction. &
\textbf{Spécification insatisfiable} & \citet{rozier2007} ; le test de satisfiabilité est le
contrôle de bon sens standard \\
\addlinespace
Un docstring affirmant « une grille dense (60\,000 points, log) » quand le code, après
corrections successives du corps, en employait 240\,001. &
\multirow{3}{3.4cm}{\textbf{Dérive de spécification} / incohérence code--commentaire} &
\multirow{3}{5.7cm}{\citet{panthaplackel2021}, et un sous-domaine de détection actif} \\
Un outil de vérification gelé avec une définition de région du ciel différant de $10^\circ$
en ascension droite de celle du corpus. & & \\
 & & \\
\addlinespace
Une validation exigeant un facteur de conversion constant à $2\%$ ; elle a passé à
$0{,}010\%$. La grandeur qu'elle convertissait n'existait pas. &
\textbf{Vacuité} & \citet{beer1997,kupferman2003} ; vingt-cinq ans d'algorithmes de détection \\
\addlinespace
Deux seuils comparés par \texttt{>=} sur une arithmétique de grille flottante ; valeurs vraies
exactement $2$ et $1$, valeurs calculées $1{,}9999999999999791$ et $0{,}9999999999999836$. Les
deux ont échoué, les deux dans le sens qui arrangeait l'analyste. &
\textbf{Test instable}, égalité flottante & \citet{luo2014} ; un lint standard \\
\addlinespace
Une validation exigeant qu'une incertitude rétrécisse d'un facteur $k$ ; la grille ne pouvait
pas résoudre une incertitude sous un pas, donc l'exigence reculait aussi vite que la grandeur
mesurée. &
\textbf{Sensibilité de test insuffisante} & la vérification de code emploie l'ordre de
précision sous raffinement de grille, pas une tolérance à grille unique
\citep{salari2000,roache1998} \\
\bottomrule
\end{tabular}
\caption{Onze instances observées dans une campagne, cartographiées sur sept classes nommées.
La compression de onze à sept est délibérée et nous la faisons nous-mêmes : les présenter
comme onze \emph{modes} distincts serait les travestir.}
\label{tab:mapping}
\end{table}

La table~\ref{tab:mapping} est l'article. Trois remarques en découlent.

\paragraph{La vacuité est celle qui devrait le plus inquiéter le praticien.} \citet{beer1997}
définissent la vacuité comme une spécification « satisfaite trop facilement, ou vacûment », et
avertissent qu'elle « induit les utilisateurs de la vérification de modèles à croire qu'un
système est correct ». Leur exemple canonique --- « toute requête est finalement suivie d'un
acquittement » est satisfaite vacûment par un système qui n'émet jamais de requête --- est
précisément notre instance, avec une tolérance à la place d'un opérateur temporel. Notre
validation est passée à $10^{-4}$ en valeur relative et n'a rien établi. \emph{Un critère qui
passe parfaitement en mesurant un fantôme est plus dangereux qu'un critère qui échoue}, et les
méthodes formelles le savent depuis avant l'existence du mouvement du pré-enregistrement. À
notre connaissance, c'est la première instance rapportée de vacuité dans un critère
d'acceptation \emph{scientifique} pré-enregistré ; c'est une revendication bien plus petite
qu'un mode nouveau, et c'est la vraie.

\paragraph{L'outillage existant aurait attrapé quatre des sept classes.} La dérive de
spécification est détectable par les outils d'incohérence code--commentaire ;
l'insatisfiabilité par un test de satisfiabilité sur le bloc de critères ; l'instabilité
flottante par un lint ; la sensibilité insuffisante par la discipline d'ordre de précision de
la vérification de code. La vacuité a des algorithmes dédiés. Rien de cet outillage n'est
appliqué aujourd'hui aux critères scientifiques pré-enregistrés, et nous ne voyons aucune
raison qu'il ne puisse pas l'être.

\paragraph{Ce qu'aucun outillage n'aurait attrapé.} La non-identifiabilité, le minimum de bord
et le comparant faible sont statistiques, non syntaxiques. Le minimum de bord en particulier
n'a été trouvé que parce qu'un lecteur humain a remarqué qu'un résultat semblait \emph{trop
beau pour être vrai} ; tous les contrôles automatiques l'avaient laissé passer.

\section{En quoi ceci diffère des taxonomies existantes}
\label{sec:yamada}

\citet{yamada2018} donne quatre façons de « craquer » le pré-enregistrement : le signalement
sélectif, la ré-expérimentation infinie, la sur-émission, et le pré-enregistrement après
connaissance des résultats. \citet{bakker2020} énumèrent 29 degrés de liberté du chercheur et
constatent que les formats structurés les réduisent sans les éliminer ; \citet{claesen2021}
auditent l'adhésion dans la première génération d'études pré-enregistrées ;
\citet{wicherts2016} donnent une liste de 34 points.

Ce sont des taxonomies \emph{comportementales} : elles classent ce que fait un chercheur, et
leurs modes sont, selon le cadrage de Yamada lui-même, ce qu'un chercheur « malveillant »
pourrait faire. Les défaillances que nous rapportons sont \emph{structurelles}. Elles classent
ce que le critère \emph{est}, et chacune s'est produite sans malhonnêteté : chacune a été
écrite par un analyste cherchant la rigueur, chacune a passé une relecture au moment de son
gel, et plusieurs ont été attrapées par la machinerie de l'analyste lui-même des jours plus
tard. Cette distinction est l'écart que nous revendiquons, et c'est le seul.

Une contiguïté n'est pas fortuite. La « ré-expérimentation infinie » de Yamada --- répéter une
expérience jusqu'à ce que les données coopèrent --- est la sœur comportementale du
\S\ref{sec:serie}, où ce sont les conceptions et non les jeux de données qui sont répétées
jusqu'à ce que l'une passe. La ressemblance de famille est exacte ; l'objet diffère.

\section{La défaillance qui n'est pas celle d'un critère}
\label{sec:serie}

Une question de la campagne a été tentée trois fois. La première conception a bloqué sur un
plancher de grille. La deuxième, à maillage dix fois plus fin, a bloqué parce que le plancher
avait été \emph{déplacé} et non levé. La troisième a remplacé l'estimateur par marche de
grille par un ajustement parabolique auto-validé --- qui a fonctionné, sans aucun plancher,
son incertitude suivant $\times2{,}00$ puis $\times4{,}00$ exactement --- et a bloqué parce que
sa propre validation exigeait la reproduction exacte des valeurs \emph{de grille} que le nouvel
estimateur abandonnait délibérément. Ce dernier blocage relève de la classe
« insatisfiabilité », et nous l'avons écrit après avoir catalogué cette classe.

Une quatrième conception tolérant un pas de grille serait passée ; toutes les autres
validations passaient déjà, et la réponse qu'elle aurait rendue était visible dans la sortie de
la troisième. Nous ne l'avons pas écrite.

\textbf{Le gel par empreinte protège chaque étude prise isolément et ne protège pas la série.}
Ce qui est gelé, c'est chaque tentative, pas la décision de recommencer. Itérer les conceptions
jusqu'à ce qu'une passe est du p-hacking d'un cran au-dessus, et le pré-enregistrement tel
qu'il se pratique n'a aucune défense contre cela --- chaque tentative est impeccablement
pré-enregistrée. Le seul garde-fou que nous voyions est une règle d'arrêt déclarée. Nous en
avons adopté une, \emph{au troisième blocage sur une même question, on cesse et on consigne le
statut non établi} ; nous notons qu'elle est arbitraire, que toute sa valeur tient à sa
publicité, et que nous ne l'avons adoptée qu'après en avoir eu besoin. Le coût est réel : une
signature restée strictement invariante sous une incertitude quatre fois moindre demeure
inopposable. Nous préférons ce coût à un verdict obtenu par sélection de conception.

\section{Pratiques adoptées, et leur antériorité}
\label{sec:pratiques}

Nous les listons par complétude et attribuons chacune ; aucune n'est de nous.

\emph{Tester le vérificateur par mutation avant de croire son taux de réussite.} Un outil
rendant « 55 sur 55, aucun échec bloquant » n'établit rien tant qu'on n'a pas montré qu'il peut
échouer ; nous avons injecté cinq fautes en exigeant que les cinq soient attrapées. C'est le
test par mutation \citep{jia2010} appliqué à l'oracle plutôt qu'au code, ce qui est
\citet{knauth2009} ; et muter les vérificateurs spécifiquement pour détecter la vacuité, c'est
\citet{date2010} --- autrement dit la combinaison de notre classe « vacuité » avec cette
pratique a été publiée en 2010. Elle a aussi été appliquée au logiciel scientifique
\citep{hook2009}. Nous rapportons seulement l'avoir trouvée nécessaire.

\emph{Détecter génériquement, jamais par liste écrite à la main.} Un contrôle de gabarit
énumérant cinq noms de marqueurs à la main est passé pendant qu'un sixième, ajouté plus tard,
restait non substitué et cassait la page publiée dès sa première ligne. Le contrôle qui l'a
remplacé est une expression régulière unique et ne peut pas manquer un marqueur dont il n'a
jamais entendu parler.

\emph{Brancher la réfutation avant la confirmation}, et rendre le seuil de confirmation plus
strict que celui de réfutation. Deux fois cet ordre a compté ; une fois il a rendu la
réfutation d'un mécanisme que nous avions affirmé dans deux entrées de registre successives
sans jamais l'avoir testé.

\emph{Déclarer le sens de chaque substitution.} Là où une valeur doit être supposée, la choisir
contre la thèse défendue, et dire que c'est pour cela.

\section{Limites}

$N=1$ : un corpus, un analyste, deux jours. Nous ne pouvons rien dire de la distribution
générale des onze instances, seulement qu'elles se sont produites et qu'aucune n'a exigé de
mauvaise foi. Le décompte est le contenu empirique : 89 critères gelés, onze défaillances des
espèces ci-dessus, et cinq études arrêtées sur leurs propres validations sans rien publier.

Le protocole est aussi mécanique. Il ne lit pas les équations et ne juge pas la physique. Dans
cette campagne, l'audit le plus conséquent a été déclenché par un lecteur humain, et une
vérification de littérature a tué net un article compagnon en établissant que ses quatre
résultats revendiqués étaient déjà publiés --- le plus ancien en 2015, le plus récent quatre
jours plus tôt. Ni l'un ni l'autre n'est atteignable par un hachage. C'est une vérification du
même genre qui a produit la forme finale de cet article : elle a établi qu'aucune de nos onze
défaillances n'était un mode nouveau, et l'article s'en trouve meilleur.

\begin{thebibliography}{99}
\bibitem[Andrae et al.(2010)]{andrae2010} Andrae, R., Schulze-Hartung, T., \& Melchior, P.
  2010, \texttt{arXiv:1012.3754}
\bibitem[Bakker et al.(2020)]{bakker2020} Bakker, M., et al. 2020, PLOS Biology, 18,
  e3000937, \texttt{doi:10.1371/journal.pbio.3000937}
\bibitem[Beaulieu-Jones \& Greene(2017)]{bj2017} Beaulieu-Jones, B.~K., \& Greene, C.~S.
  2017, Nature Biotechnology, 35, 342, \texttt{doi:10.1038/nbt.3780}
\bibitem[Beer et al.(1997)]{beer1997} Beer, I., Ben-David, S., Eisner, C., \& Rodeh, Y.
  1997, in CAV 1997, \texttt{doi:10.1007/3-540-63166-6\_28}
\bibitem[Claesen et al.(2021)]{claesen2021} Claesen, A., et al. 2021, Royal Society Open
  Science, 8, 211037, \texttt{doi:10.1098/rsos.211037}
\bibitem[Ferrari Dacrema et al.(2019)]{dacrema2019} Ferrari Dacrema, M., Cremonesi, P., \&
  Jannach, D. 2019, in RecSys 2019, \texttt{doi:10.1145/3298689.3347058}
\bibitem[Feldman \& Cousins(1998)]{feldman1998} Feldman, G.~J., \& Cousins, R.~D. 1998,
  Phys.\ Rev.\ D, 57, 3873, \texttt{doi:10.1103/PhysRevD.57.3873}
\bibitem[Gipp et al.(2017)]{cryptsubmit} Gipp, B., Breitinger, C., Meuschke, N., \&
  Beel, J. 2017, in JCDL 2017, \texttt{doi:10.1109/JCDL.2017.7991588}
\bibitem[Gipp et al.(2018)]{originstamp} Gipp, B., et al. 2018, it --- Information
  Technology, 60, 47, \texttt{doi:10.1515/itit-2018-0020}
\bibitem[Hook \& Kelly(2009)]{hook2009} Hook, D., \& Kelly, D. 2009, in ICSE SECSE 2009,
  \texttt{doi:10.1109/SECSE.2009.5069163}
\bibitem[Jia \& Harman(2010)]{jia2010} Jia, Y., \& Harman, M. 2010, IEEE Trans.\ Softw.\
  Eng., 37, 649, \texttt{doi:10.1109/TSE.2010.62}
\bibitem[Jimenez et al.(2017)]{popper} Jimenez, I., et al. 2017, in IPDPSW 2017,
  \texttt{doi:10.1109/IPDPSW.2017.157}
\bibitem[Klein \& Roodman(2005)]{klein2005} Klein, J.~R., \& Roodman, A. 2005, Ann.\ Rev.\
  Nucl.\ Part.\ Sci., 55, 141, \texttt{doi:10.1146/annurev.nucl.55.090704.151521}
\bibitem[Knauth et al.(2009)]{knauth2009} Knauth, T., Fetzer, C., \& Felber, P. 2009, in
  ICSTW 2009, \texttt{doi:10.1109/ICSTW.2009.40}
\bibitem[Kupferman \& Vardi(2003)]{kupferman2003} Kupferman, O., \& Vardi, M.~Y. 2003,
  STTT, 4, 224, \texttt{doi:10.1007/s100090100062}
\bibitem[Luo et al.(2014)]{luo2014} Luo, Q., Hariri, F., Eloussi, L., \& Marinov, D. 2014,
  in FSE 2014, \texttt{doi:10.1145/2635868.2635920}
\bibitem[MacCoun \& Perlmutter(2015)]{maccoun2015} MacCoun, R., \& Perlmutter, S. 2015,
  Nature, 526, 187, \texttt{doi:10.1038/526187a}
\bibitem[Muir et al.(2020)]{muir2020} Muir, J., et al. 2020, MNRAS, 494, 4454,
  \texttt{doi:10.1093/mnras/staa965}
\bibitem[Panthaplackel et al.(2021)]{panthaplackel2021} Panthaplackel, S., et al. 2021, in
  AAAI 2021, \texttt{doi:10.1609/aaai.v35i1.16119}
\bibitem[Roache(1998)]{roache1998} Roache, P.~J. 1998, AIAA Journal, 36, 696,
  \texttt{doi:10.2514/2.457}
\bibitem[Rozier \& Vardi(2007)]{rozier2007} Rozier, K.~Y., \& Vardi, M.~Y. 2007, in SPIN
  2007, \texttt{doi:10.1007/978-3-540-73370-6\_11}
\bibitem[Salari \& Knupp(2000)]{salari2000} Salari, K., \& Knupp, P. 2000, rapport Sandia
  SAND2000-1444, \texttt{doi:10.2172/759450}
\bibitem[Vacuité par mutation(2010)]{date2010} 2010, « Vacuity analysis for property
  qualification by mutation of checkers », in DATE 2010,
  \texttt{doi:10.1109/DATE.2010.5457158}
\bibitem[Wicherts et al.(2016)]{wicherts2016} Wicherts, J.~M., et al. 2016, Frontiers in
  Psychology, 7, 1832, \texttt{doi:10.3389/fpsyg.2016.01832}
\bibitem[Willroth \& Atherton(2024)]{willroth2024} Willroth, E.~C., \& Atherton, O.~E.
  2024, AMPPS, 7, \texttt{doi:10.1177/25152459231213802}
\bibitem[Yamada(2018)]{yamada2018} Yamada, Y. 2018, Frontiers in Psychology, 9, 1831,
  \texttt{doi:10.3389/fpsyg.2018.01831}
\end{thebibliography}

\end{document}
